% Part: first-order-logic
% Chapter: introduction
% Section: first-order-logic

\documentclass[../../../include/open-logic-section]{subfiles}

\begin{document}

\olfileid[pt]{fol}{int}{fol}

\olsection{Lógica de Primeira Ordem}

Você provavelmente já conhece a lógica de primeira ordem da sua
primeira introdução à lógica formal.\footnote{De fato, mais ou menos
assumimos que sim!{} Se não conhece, pode revisar um livro-texto mais
elementar, como \emph{forall x} \citep{Magnus2021}.} Pode conhecê-la
como ``lógica quantificacional'' ou ``lógica de predicados''.  A
lógica de primeira ordem, antes de tudo, é uma linguagem formal.  Isso
significa que ela tem um certo vocabulário, e suas expressões são
cadeias desse vocabulário.  Mas nem toda cadeia é permitida.  Há
diferentes tipos de expressões permitidas: termos, !!{formula}s e
!!{sentence}s.  Estamos principalmente interessados em !!{sentence}s da
lógica de primeira ordem: elas nos fornecem um análogo formal de
sentenças do inglês, e sobre elas podemos fazer as perguntas que um
lógico tipicamente se interessa. Por exemplo:
\begin{itemize}
    \item $!B$ segue logicamente de~$!A$?
    \item $!A$ é logicamente verdadeira, logicamente falsa ou
contingente?
    \item $!A$ e $!B$ são equivalentes?
\end{itemize}

Essas perguntas são primariamente perguntas sobre o ``significado'' das
!!{sentence}s da lógica de primeira ordem.  Por exemplo, um filósofo
analisaria a pergunta de se $!B$ segue logicamente de~$!A$ como
perguntando: há um caso em que $!A$ é verdadeira mas~$!B$ é falsa
($!B$ não segue de~$!A$), ou todo caso que torna $!A$ verdadeira
também torna~$!B$ verdadeira ($!B$ segue de~$!A$)?  Mas ainda não nos
foi dito o que é um ``caso''---esse é o trabalho da \emph{semântica}.
A semântica da lógica de primeira ordem fornece um modelo matematicamente
preciso da ideia intuitiva do filósofo de ``caso'', e também---e isso
é importante---do que significa uma !!{sentence}~$!A$ ser \emph{verdadeira
em} um caso. Chamamos o modelo matematicamente preciso que
desenvolveremos de !!a{structure}. A relação que torna ``verdadeira em''
precisa é chamada de relação de \emph{satisfação}.  Assim, o que
definiremos é ``$!A$ é satisfeita em~$\Struct{M}$'' (em símbolos:
$\Sat{M}{!A}$) para !!{sentence}s~$!A$ e !!{structure}s~$\Struct{M}$. Uma vez
feito isso, também podemos dar definições precisas dos outros termos
semânticos como ``segue de'' ou ``é logicamente verdadeira''. Essas
definições tornarão possível decidir, novamente com precisão
matemática, se, por exemplo, $\lforall[x][(!A(x) \lif !B(x)),
\lexists[x][!A(x)] \Entails \lexists[x][!B(x)]]$. A resposta será, é
claro, ``sim''. Se você já foi treinado para simbolizar sentenças do
inglês na lógica de primeira ordem, reconhecerá isto como, por
exemplo, a simbolização de ``Todas as formigas são insetos, há
formigas, portanto há insetos.'' Isso é obviamente um argumento
válido, e assim nosso modelo matemático de ``segue de'' para nossa
linguagem formal deve dar a mesma resposta.

Outro tópico que você provavelmente lembra da sua primeira introdução
à lógica formal é que há \emph{!!{derivation}s}.  Se você fez um
primeiro curso de lógica formal, seu professor o fez praticar a
encontrar tais !!{derivation}s, talvez até uma !!{derivation} que
mostra que a consequência acima vale.  Há muitas maneiras diferentes
de dar !!{derivation}s: você pode ter feito algo chamado ``dedução
natural'' ou ``árvores de verdade'', mas há muitas outras. O propósito
dos sistemas de !!{derivation} é fornecer ferramentas com as quais as
perguntas dos lógicos acima podem ser respondidas: por exemplo, uma
!!{derivation} de dedução natural em que $\lforall[x][(!A(x) \lif !B(x))]$
e $\lexists[x][!A(x)]$ são premissas e $\lexists[x][!B(x)]$ é a
conclusão (última linha) \emph{verifica} que $\lexists[x][!B(x)]$
segue logicamente de $\lforall[x][(!A(x) \lif !B(x))]$ e
$\lexists[x][!A(x)]$.  

Mas por quê?  À primeira vista, os sistemas de !!{derivation} não têm
nada a ver com semântica: dar uma !!{derivation} formal envolve
apenas dispor símbolos de certas maneiras regidas por regras; eles não
mencionam ``casos'' ou ``verdadeira em'' de forma alguma.  A conexão
entre sistemas de !!{derivation} e semântica tem de ser estabelecida por
uma investigação metalógica. O que é necessário é uma prova
matemática, por exemplo, de que uma !!{derivation} formal de
$\lexists[x][!B(x)]$ a partir das premissas $\lforall[x][(!A(x) \lif
!B(x))]$ e $\lexists[x][!A(x)]$ é possível se, e somente se,
$\lexists[x][!B(x)]$ é consequência de $\lforall[x][(!A(x) \lif !B(x))]$
e $\lexists[x][!A(x)]$ em conjunto.  Antes que isso possa ser feito,
entretanto, muito trabalho minucioso tem de ser realizado para acertar
as definições de sintaxe e semântica.

\end{document}
